% =====================================================================
% Example Section
% ---------------------------------------------------------------------
% Every cross-reference in this template uses cleveref, in IEEE style:
%   \cref{key}  -> "Fig. 1", "Fig. 2(a)", "Table I", "(3)", "Section II-A"  (mid-sentence)
%   \Cref{key}  -> same, except equations: "Equation (3)"   (start of a sentence)
%   \cref{a,b}  -> "Figs. 1 and 2"   (lists are sorted and compressed)
%   \crefrange{a}{c} -> "Figs. 1--3"
% Label with a prefix so keys are easy to find: sec:, fig:, tab:, eq:, lst:, app:
% Never type "Fig.~\ref{...}" by hand; cleveref adds the name for you.
%
% Every figure and table uses [H] (from the float package) so it is placed
% exactly where it appears in the source instead of floating elsewhere.
% =====================================================================

\section{Example Section}
\label{sec:example}

\subsection{Figures}

\lipsum[2]

A single figure is shown in \cref{fig:single}. Keep your images in the
\texttt{Figures/} folder and include them by path; the file extension can
be left off. JPG and PNG suit photos and screenshots, while PDF keeps
plots and diagrams exported from MATLAB, KiCad and similar programs sharp
at any zoom. Always place the \verb|\label| \emph{after} the
\verb|\caption|, otherwise the reference will point to the section instead
of the figure. IEEE style also avoids starting a caption with ``A'',
``An'' or ``The''.

\begin{figure}[H]
    \centering
    \includegraphics[width=0.6\linewidth]{Figures/test-setup}
    \caption{Example photograph loaded from the \texttt{Figures/} folder.}
    \label{fig:single}
\end{figure}

\subsubsection{Side-by-side subfigures}

\Cref{fig:pair} contains two subfigures. Each subfigure can be referenced
on its own, e.g.\ \cref{fig:pair-a} shows the input and \cref{fig:pair-b}
shows the response. Several labels can be passed at once:
\cref{fig:single,fig:pair}.

In IEEE style, each subfigure is marked only with its letter, using an
empty \verb|\caption{}| followed by its \verb|\label|. The description of
each part goes in the main caption after the figure title, e.g.\
\verb|\caption{Title. (a) First part. (b) Second part.}|.

\begin{figure}[H]
    \centering
    \begin{subfigure}[b]{0.48\linewidth}
        \centering
        \includegraphics[width=\linewidth]{example-image-a}
        \caption{}
        \label{fig:pair-a}
    \end{subfigure}
    \hfill
    \begin{subfigure}[b]{0.48\linewidth}
        \centering
        \includegraphics[width=\linewidth]{example-image-b}
        \caption{}
        \label{fig:pair-b}
    \end{subfigure}
    \caption{Two subfigures placed side by side. (a) Step input. (b) System response.}
    \label{fig:pair}
\end{figure}

\subsubsection{Grid of subfigures}

A $2 \times 2$ grid is built by inserting a blank line (paragraph break)
after every second subfigure, as in \cref{fig:grid}. The four panels,
\crefrange{fig:grid-a}{fig:grid-d}, could be different test conditions.

\begin{figure}[H]
    \centering
    \begin{subfigure}[b]{0.45\linewidth}
        \centering
        \includegraphics[width=\linewidth]{example-image-a}
        \caption{}
        \label{fig:grid-a}
    \end{subfigure}
    \hfill
    \begin{subfigure}[b]{0.45\linewidth}
        \centering
        \includegraphics[width=\linewidth]{example-image-b}
        \caption{}
        \label{fig:grid-b}
    \end{subfigure}

    \vspace{1em}

    \begin{subfigure}[b]{0.45\linewidth}
        \centering
        \includegraphics[width=\linewidth]{example-image-c}
        \caption{}
        \label{fig:grid-c}
    \end{subfigure}
    \hfill
    \begin{subfigure}[b]{0.45\linewidth}
        \centering
        \includegraphics[width=\linewidth]{example-image}
        \caption{}
        \label{fig:grid-d}
    \end{subfigure}
    \caption{$2 \times 2$ grid of subfigures. (a) \qty{25}{\celsius}. (b) \qty{50}{\celsius}. (c) \qty{75}{\celsius}. (d) \qty{100}{\celsius}.}
    \label{fig:grid}
\end{figure}

\subsection{Tables}

All tables use \texttt{booktabs}: \verb|\toprule|, \verb|\midrule| and
\verb|\bottomrule| replace \verb|\hline|, and vertical lines are not used.
Table captions go \emph{above} the table and, in IEEE style, do not end
with a full stop.

\subsubsection{Basic table}

\Cref{tab:basic} is the simplest booktabs layout.

\begin{table}[H]
    \centering
    \caption{Material properties at room temperature}
    \label{tab:basic}
    \begin{tabular}{lccc}
        \toprule
        Material  & Density (\unit{\kilo\gram\per\cubic\metre}) & $E$ (\unit{\giga\pascal}) & $\nu$ \\
        \midrule
        Steel     & 7850 & 200 & 0.30 \\
        Aluminium & 2700 & 69  & 0.33 \\
        Titanium  & 4500 & 116 & 0.34 \\
        \bottomrule
    \end{tabular}
\end{table}

\subsubsection{Numbers aligned on the decimal point}

\Cref{tab:siunitx} uses \texttt{siunitx} \texttt{S} columns so values line
up on the decimal point. \verb|table-format=3.2| means up to 3 digits before
and 2 after the decimal point. Wrap text headers in braces \verb|{...}|.

\begin{table}[H]
    \centering
    \caption{Measured voltage and current for the test circuit}
    \label{tab:siunitx}
    \begin{tabular}{
        l
        S[table-format=2.2]
        S[table-format=1.3]
        S[table-format=3.1]
    }
        \toprule
        {Test} & {Voltage (\unit{\volt})} & {Current (\unit{\ampere})} & {Power (\unit{\watt})} \\
        \midrule
        T1 & 5.00  & 0.125 & 0.6   \\
        T2 & 12.00 & 1.040 & 12.5  \\
        T3 & 24.50 & 4.210 & 103.1 \\
        \bottomrule
    \end{tabular}
\end{table}

\subsubsection{Grouped columns}

\Cref{tab:grouped} groups columns with \verb|\multicolumn| and underlines
each group with \verb|\cmidrule(lr){a-b}|. The \verb|(lr)| trims the rule
so neighbouring groups do not touch. \verb|\addlinespace| separates row
groups without drawing a line.

\begin{table}[H]
    \centering
    \caption{Comparison of simulated and experimental results}
    \label{tab:grouped}
    \begin{tabular}{l S[table-format=2.1] S[table-format=2.1] S[table-format=2.1] S[table-format=2.1]}
        \toprule
        & \multicolumn{2}{c}{Simulation} & \multicolumn{2}{c}{Experiment} \\
        \cmidrule(lr){2-3} \cmidrule(lr){4-5}
        {Case} & {Rise (\unit{\milli\second})} & {Overshoot (\%)} & {Rise (\unit{\milli\second})} & {Overshoot (\%)} \\
        \midrule
        P controller   & 12.4 & 18.2 & 13.1 & 20.5 \\
        PI controller  & 10.8 & 9.6  & 11.5 & 11.2 \\
        \addlinespace
        PID controller & 8.2  & 4.1  & 9.0  & 5.3  \\
        \bottomrule
    \end{tabular}
\end{table}

\subsection{Equations and Code}

Equations are referenced the same way. \Cref{eq:ohm} is Ohm's law and
\cref{eq:power} gives the power used to fill in \cref{tab:siunitx}.
\begin{equation}
    V = IR
    \label{eq:ohm}
\end{equation}
\begin{equation}
    P = VI = I^2 R
    \label{eq:power}
\end{equation}

Code listings take the label as an option, as in \cref{lst:example}. Longer
listings belong in the appendix, see \cref{app:code}.

\begin{lstlisting}[caption={Calculating power from measured data.}, label={lst:example}]
V = [5.00 12.00 24.50];     % Voltage (V)
I = [0.125 1.040 4.210];    % Current (A)
P = V .* I;                 % Power (W)
\end{lstlisting}

\cleardoublepage
